\documentclass[10pt,a4paper]{amsart}
\usepackage[margin=19mm]{geometry}
\usepackage{amsmath,amssymb,mathtools}
\usepackage[scr=rsfs,cal=euler]{mathalfa}
\usepackage{xfrac}
\usepackage[unicode,hidelinks,bookmarks=false]{hyperref}
\newcommand{\ie}{i.e.\ }
\newcommand{\cf}{cf.\ }
\newcommand{\resp}{respectively\ }
\newcommand{\Subsection}{\subsection}
\providecommand{\nobreakdash}{\nobreak\hbox{-}}
\setlength{\emergencystretch}{2em}
\multlinegap0pt
\usepackage{tikz}
\usepackage{latexsym,epic}
\usepackage{latexsym,epic}
\usepackage{color}
\usepackage{soul}
\hypersetup{
	colorlinks=true, %set true if you want colored links
	linktoc=all, %set to all if you want both sections and subsections linked
	linkcolor=blue} %choose some color if you want links to stand out
\newtheorem{theorem}{Theorem}[section]
\newtheorem{lemma}[theorem]{Lemma}
\newtheorem{corollary}[theorem]{Corollary}
\newtheorem{proposition}[theorem]{Proposition}
\newtheorem{definition}[theorem]{Definition}
\newtheorem{conjecture}[theorem]{Conjecture}
\newtheorem{example}[theorem]{Example}
\newtheorem{remark}[theorem]{Remark}
\begin{document}
\title[Arithmetic Properties Satisfied by a Recent Integer Partitio]{Arithmetic Properties Satisfied by a Recent Integer Partition Function of Dombos}
\author{Robson da Silva; James A. Sellers}
\date{}
\maketitle
\noindent\emph{Source and licence.} Arithmetic Properties Satisfied by a Recent Integer Partition Function of Dombos. Version arXiv:2606.17440v1 (CC BY 4.0). This material is distributed under the Creative Commons Attribution 4.0 International licence, http://creativecommons.org/licenses/by/4.0/, as recorded in the arXiv OAI licence metadata for this exact version and shown on the abstract page https://arxiv.org/abs/2606.17440v1. Full rights notice: licenses/arxiv-2606.17440-CC-BY-4.0.txt.\par\smallskip
\noindent\textit{Editorial scope.} Counted unit: the original \texttt{lemma} environment \texttt{gf2} at source lines 157--162 together with its complete proof at lines 164--184; the delivered document keeps source lines 89--185 so that every referenced label and every citation resolves inside it. Native editable source is the paper's own concatenated TeX; the AMS wrapper is editorial formatting only. No publisher class, upstream script or unrelated artwork is redistributed.
\begin{equation}
G(q):=\sum_{n=0}^{\infty} dp(n)q^n = \frac{1}{(q;q^6)_\infty (q^5;q^6)_\infty (q^4;q^4)_\infty},
\label{gf1}
\end{equation}
where $(a;q)_\infty := (1-a)(1-aq)(1-aq^{2})\cdots$
is the usual $q$-Pochhammer symbol. 

The primary goal of this brief note is to prove a number of Ramanujan--like congruences satisfied by the function $dp(n)$ in certain arithmetic progressions.  In the work below, we will prove the following results. 

\begin{theorem}
For all $n \geq 0$,
	\begin{align}
	dp(6n+4) & \equiv 0 \pmod{2}, \label{c1} \\ 
%    dp(27n+16) & \equiv 0 \pmod{3}, \label{c2} \\
    dp(18n+10) & \equiv 0 \pmod{4}, \label{c3} \\
    dp(54n+52) &  \equiv 0 \pmod{8}. \label{c4}
    \end{align}
\label{Th1}
\end{theorem}

\begin{theorem} Let $p$ be a prime such that $p \equiv 17 \mbox{ or } 23 \pmod{24}$. Then for all $k,m \geq 0$ with $p \nmid m$, we have
\begin{equation*}
dp\left( 6p^{2k+1}m + \frac{15p^{2k+2}+1}{4} \right) \equiv 0 \pmod{4}.
\end{equation*}
\label{Th2}
\end{theorem}

\begin{theorem}
For all $n \geq 0$,
\begin{align}
dp(27n+16) & \equiv 0 \pmod{3}.
\label{c2}
\end{align}
\label{Th3}
\end{theorem}

\begin{theorem}
For all $n \geq 0$,
\begin{equation*}
dp(27n+7) \equiv dp(3n+1) \pmod{3}.
\end{equation*}
\label{internal_cong}
\end{theorem}

\begin{theorem}
For all $\alpha \geq 1$ and all $n \geq  0$,
\begin{equation*}
    dp \left( 3^{2\alpha + 1}n + \frac{7 \cdot 9^\alpha + 1}{4} \right) \equiv 0 \pmod{3}.
\end{equation*}
\label{Th4}
\end{theorem}

	







\section{Preliminaries}

Throughout the remainder of this paper, we define 
		$$f_k := (q^k;q^k)_{\infty}$$
in order to shorten the notation.

We begin by rewriting the generating function \eqref{gf1} in the following form.


\begin{lemma}
\begin{equation}
\sum_{n=0}^{\infty} dp(n)q^n = \frac{f_2f_3}{f_1f_4f_6}.
\label{gf2}
\end{equation}
\end{lemma}

\begin{proof}
Initially, we note that
\begin{align*}
\frac{1}{(q;q^6)_\infty (q^5;q^6)_\infty} & = \frac{(q^2;q^6)_\infty(q^3;q^6)_\infty(q^4;q^6)_\infty(q^6;q^6)_\infty}{(q;q^6)_\infty(q^2;q^6)_\infty(q^3;q^6)_\infty(q^4;q^6)_\infty(q^5;q^6)_\infty(q^6;q^6)_\infty} \\
& = \frac{(q^3;q^3)_\infty(q^2;q^6)_\infty(q^4;q^6)_\infty}{(q;q^3)_\infty(q^2;q^3)_\infty(q^3;q^3)_\infty} \\
& = \frac{(q^2;q^6)_\infty(q^4;q^6)_\infty}{(q;q^3)_\infty(q^2;q^3)_\infty} \\
& = \frac{(q^2;q^2)_\infty}{(q^6;q^6)_\infty}\cdot \frac{(q^3;q^3)_\infty}{(q;q)_\infty} \\
& = \frac{f_2f_3}{f_1f_6}.  
%& = \frac{(-q;q^3)_\infty(-q^2;q^3)_\infty(q;q^3)_\infty(q^2;q^3)_\infty}{(q;q^3)_\infty(q^2;q^3)_\infty} \\
%& = (-q;q^3)_\infty(-q^2;q^3)_\infty.
\end{align*}
Thus, it follows from \eqref{gf1} that
\begin{align*}
 \sum_{n=0}^{\infty} dp(n)q^n 
 %& = \frac{1}{f_4} (-q;q^3)_\infty(-q^2;q^3)_\infty \\
 %& = \frac{1}{f_4} \frac{(-q;q^3)_\infty(-q^2;q^3)_\infty (-q^3;q^3)_\infty}{(-q^3;q^3)_\infty} \\
 %& = \frac{1}{f_4} \frac{(-q;q)_\infty}{(-q^3;q^3)_\infty} \\
 & = \frac{f_2f_3}{f_1f_4f_6},
\end{align*}
which is \eqref{gf2}.
\end{proof}


\end{document}
